\subsection{Итоговые проекты курса}

Поздравляем! Вы освоили управление дроном Pioneer на языке Python. Теперь пришло время применить знания в комплексных проектах.

\subsubsection*{Проект 1: Автономный инспектор (Склад)}
\textbf{Задача:} Дрон должен облететь 4 стеллажа (маркеры ID 0, 1, 2, 3), расположенных по углам комнаты.
\textbf{Алгоритм:}
1. Взлет.
2. Поиск маркера 0. Подлет к нему на 1 метр.
3. Поворот на 90 градусов, поиск маркера 1.
4. ...
5. Возврат на точку старта и посадка.

\subsubsection*{Проект 2: <<Следуй за мной>> (Follow Me)}
\textbf{Задача:} Дрон должен удерживать человека в центре кадра на безопасном расстоянии.
\textbf{Требования:}
1. Использовать MediaPipe для детекции человека.
2. Если человек подходит близко (плечи занимают весь кадр) — дрон отлетает назад.
3. Если человек уходит — дрон летит вперед.
4. Поворот за смещением человека влево/вправо.

\subsubsection*{Проект 3: Гонка по точкам}
\textbf{Задача:} Пролететь через полосу препятствий (ворота), отмеченные ArUco-маркерами, как можно быстрее.
\textbf{Усложнение:} Маркеры могут быть расположены на разной высоте. Дрон должен пролетать сквозь ворота (или над ними).

\subsubsection*{Критерии оценки проектов}
\begin{enumerate}
    \item \textbf{Стабильность:} Программа не падает с ошибками.
    \item \textbf{Безопасность:} Дрон не совершает резких опасных маневров.
    \item \textbf{Читаемость кода:} Код разбит на функции, есть комментарии.
    \item \textbf{Функциональность:} Задача выполнена полностью.
\end{enumerate}
